% Nouveau EGA V (cote 25): archive cote 25, site page 116 (PDF page 117).
% typescript. Draft compared: none.
% AI-transcribed by Claude Opus 5.5 (claude-opus-5-5, adaptive thinking, effort high), Message Batches API, batch ega25b (msgbatch_01A7n5YuXubUPNK1zsepi8RG), 2026-11-13; checked against both supplied scan parts; see CHECK_DECISIONS.tsv.
\scanpage{25/116}
\mspage{19}

points fixes, elle est déterminée \xout{uniquement} par le morphisme f:X $\hand{\longrightarrow}$ P
associé. Noter d'ailleurs qu'on \deleted{\xout{\illegible{1}}}obtient évidemment ainsi une correspon
dance biunivoque entre systêmes linéaires de diviseurs de X/S paramétré
par $\hand{\check{P}}$ , et morphismes f:X $\hand{\longrightarrow}$ P tel que \deleted{\xout{\illegible{1}}}$(fx_S\mathrm{id})^{-1}(\underline{H})$ soit un
diviseur\deleted{ de Cartier} \xout{sur} relatif sur $Xx_S\overset{\hand{\vee}}{P}$ par rapport à $\overset{\hand{\vee}}{P}$ . Ce\over{\unsure{ci}}{tte} \deleted{se} \inserted{condition}
\inserted{se} vérifie aussi fibre par fibre, et on obtient :

\uline{Proposition} \over{\unsure{1|2}}{3}.\hand{1}.On a une correspondance biunivoque entre les \xout{morphismes}
systêmes linéaires \inserted[typed]{sans points fixes} de diviseurs sur X/S paramétrés par $\hand{\check{P}}$ , et les morphis
mes f: X $\hand{\longrightarrow}$ P ayant la propriété suivante: pour tout $s\hand{\in} S$, d\over{e}{é}signant
par $\bar{k}$ une clôture algébrique de k(s), et pour tout cycle premier associé
X' de $X_{\bar{k}}$ , $f(X')\hand{\subset} P_{\bar{k}}$ n'est contenu dans aucun hyperplan de $P_{\bar{k}}$ .
\xout{(NB}\hand{(}Lorsque X est à fibres géométriquement \xout{\unsure{intégres|intègres}} \inserted[typed]{intègres}, ceci \xout{se simplifie}
\xout{et} s'énoncé simplement en disant que $f(X_{\bar{k}})$ n'est contenu dans aucun
hyperplan de $P_{\bar{k}}$ \hand{)}.

En général (i.e. \xout{en} lorsque Z $\hand{\neq}\hand{\emptyset}$ ) on ne peut plus affirmer que
la connaissance de f détermine \unsure{D|l}a famille de diviseurs. Le cas le plus
trivial est celui où $\hand{\check{P}}$ = S \inserted[typed]{est de dim relative 0}, alors la donnée d'un systême linéaire de\hand{.}
diviseurs de X/S paramétrée par S équivaut à celle d'un diviseur de
Cartier relatif D sur X par rapport à S; \xout{Tout l'app\illegible{1}} le morphisme
associé est la projection X-D $\hand{\longrightarrow}$ S , et on voit \xout{(même en} que la connais
sance de ce morphisme (y compris celle de son domaine de définition)
ne détermine pas D. Dans ce cas également, supposant pour simplifier X
réduit, le domaine de définition de \xout{ce morphisme} f \xout{n'est pas} considéré
comme application rationelle de X dans P = S n'est pas X-D , mais X.
Pour éliminer ce genre de phénomènes déplaisants, il y a lieu de se bor-%
ner aux systêmes \unsure{h|l}i\unsure{\'n|n}éaires de diviseurs ``\uline{sans composante fixe}'' . De
façon générale, si S=Spec(k), étaht donné une famille (non néc linéaire)

\margin{116}
